\chapter{Navier-Stokes validation: the lid-driven cavity vs Ghia (1982)}
\label{ch:D2_lid_driven_cavity}
\noindent\texttt{tutorials/D\_flow/D2\_lid\_driven\_cavity.py}\\[0.75em]
\begin{outcome}
You can validate a flow solver against archival
literature data and state the tolerance honestly (a 33x33 grid vs Ghia's
129x129). You meet the linearized monolithic stepper: ONE (u,p) solve per
pseudo-time step, convection frozen at the previous velocity.

THE PROBLEM. The classic incompressible-flow benchmark: a unit box of fluid,
three no-slip walls, and a lid sliding at u = 1. At Re = 100 the flow settles
into one primary vortex. Ghia, Ghia \& Shin (1982) tabulated the steady
centerline velocity profiles from a 129x129 fine-grid solve; every CFD code
since has been judged against those two columns of numbers. So is ours.

WHAT IS RUNNING UNDERNEATH.
- Equal-order Q1/Q1 velocity-pressure on the octree, kept stable by
  VMS/PSPG stabilization (tau\_M-weighted residual terms; grad-div tau\_C).
- Convection in the s = 1/2 skew form  a.grad(u) + (1/2) div(a) u  — the
  discretely ENERGY-STABLE choice: the advection operator is exactly
  skew-symmetric, so convection cannot create kinetic energy (we verify this
  as a test, and it is also what will make the adjoint elegant in M1c).
- One LINEAR solve per pseudo-time step: convection is frozen at the
  previous velocity (the production "linearized monolithic" stepper),
  marched with BDF1 until nothing changes.
\end{outcome}

\section*{The script}
\begin{lstlisting}
import numpy as np

from diffsim.octree.build import build_uniform
from diffsim.mesh.nodes import build_mesh
from diffsim.mesh.constraints import build_constraints
from diffsim.mesh.basis import basis_tables
from diffsim.assembly.operators import DeviceMesh
from diffsim.steppers.linearized import LinearizedMonolithicStepper
from diffsim.mesh.pointeval import point_eval_weights

DEVICE = "cuda:0"

# Ghia, Ghia & Shin (1982), Re = 100: u on the vertical centerline x = 0.5
GHIA_Y = np.array([0.0000, 0.0547, 0.0625, 0.0703, 0.1016, 0.1719, 0.2813,
                   0.4531, 0.5000, 0.6172, 0.7344, 0.8516, 0.9531, 0.9609,
                   0.9688, 0.9766, 1.0000])
GHIA_U = np.array([0.0000, -0.03717, -0.04192, -0.04775, -0.06434, -0.10150,
                   -0.15662, -0.21090, -0.20581, -0.13641, 0.00332, 0.23151,
                   0.68717, 0.73722, 0.78871, 0.84123, 1.00000])


def lid(x, t):
    g = np.zeros((len(x), 2))
    g[np.abs(x[:, 1] - 1.0) < 1e-12, 0] = 1.0     # the moving lid
    return g


if __name__ == "__main__":
    level, Re, dt = 5, 100.0, 0.05
    tree = build_uniform(level, dim=2)
    mesh = build_mesh(tree, p=1)
    cons = build_constraints(mesh)
    dm = DeviceMesh.from_mesh(mesh, cons, basis_tables(1, dim=2), DEVICE)
    stepper = LinearizedMonolithicStepper(
        dm, nu=1.0 / Re, dt=dt,
        f_fn=lambda x, t: np.zeros((len(x), 2)), g_fn=lid, order=1)
    stepper.set_initial(lambda x: np.zeros((len(x), 2)))

    prev = None
    for step in range(1, 401):
        x = stepper.step()
        u = x[:, :2]
        if prev is not None and np.abs(u - prev).max() / dt < 2e-4:
            print(f"steady after {step} pseudo-time steps")
            break
        prev = u.copy()

    W = point_eval_weights(mesh, np.stack(
        [np.full_like(GHIA_Y, 0.5), GHIA_Y], axis=1))
    u_c = np.asarray(W @ np.asarray(dm.constraints.T @ u[:, 0]))

    print(f"\n{'y':>8} {'u (ours)':>10} {'u (Ghia)':>10} {'diff':>9}")
    for y, uo, ug in zip(GHIA_Y, u_c, GHIA_U):
        print(f"{y:>8.4f} {uo:>10.4f} {ug:>10.4f} {uo - ug:>9.4f}")
    print(f"\nmax |diff| = {np.abs(u_c - GHIA_U).max():.4f}   "
          f"(a 33x33 grid vs Ghia's 129x129 — tolerance 0.06)")


# PERFORMANCE CORNER: each pseudo-time step assembles AND factorizes (the
# advecting field changes). Time assemble vs splu per step at level 6; then
# read src/diffsim/solvers/krylov_dev.py and estimate the speedup of
# replacing splu with the fused BiCGStab (P2 measures it).
\end{lstlisting}

\begin{explore}
\begin{verbatim}
(a) Refine to level 6 and watch the profile tighten (runtime ~4x).
  (b) Swap in the Leray pressure-projection stepper (see
      tests/test_cavity.py::test_cavity_re100_ghia_leray) — same physics
      through a completely different time-splitting; the profiles must
      agree. They do, to 0.0035.
  (c) Raise Re to 400. What breaks first: steadiness of the pseudo-time
      march, or the coarse-grid profile? (Ghia's Re=400 column is in the
      1982 paper.)
\end{verbatim}
\end{explore}
